\chapter{Диагностика и отладка}

Большинство проблем с Docker решается тремя вопросами: что говорит
сам контейнер, что он видит изнутри и что происходит на хосте.
В главе собраны команды и приёмы, которые позволяют быстро пройти
этот путь: от \cmd{docker ps} до \cmd{nsenter} и анализа OOM.

\section{С чего начать}

\begin{enumerate}
  \item \cmd{docker ps -a}~--- какие контейнеры есть, в каком они
        состоянии и с каким кодом выхода.
  \item \cmd{docker logs <контейнер>}~--- что писало приложение.
  \item \cmd{docker inspect <контейнер>}~--- как он сконфигурирован:
        команда, переменные, монтирования, сеть, лимиты.
  \item Проверка изнутри: \cmd{docker exec} или временный отладочный
        контейнер.
  \item Проверка хоста: диск, память, порты, журнал демона.
\end{enumerate}

\section{Инспекция контейнера}

Команда \cmd{docker inspect} возвращает большой JSON. Полезно сразу
извлекать нужные поля через \flag{--format}.

\begin{code}[language=bash]
$ docker inspect --format '{{.State.Status}}' web
running
$ docker inspect --format '{{.State.ExitCode}} {{.State.OOMKilled}}' web
137 true
$ docker inspect --format '{{json .Config.Env}}' web
["PATH=/usr/local/sbin:...","NGINX_VERSION=1.27.3"]
$ docker inspect --format '{{json .Mounts}}' web
$ docker inspect --format '{{json .NetworkSettings.Networks}}' web
$ docker inspect --format '{{.HostConfig.Memory}}' web
536870912
$ docker inspect --format '{{json .Config.Labels}}' web
\end{code}

\begin{tipbox}
Сохраните полезные шаблоны в алиасы:
\cmd{docker inspect -f '\{\{.State.Status\}\} \{\{.State.ExitCode\}\}'}.
Это быстрее, чем читать весь JSON глазами.
\end{tipbox}

\section{Логи и события}

\begin{code}[language=bash]
$ docker logs --tail 200 web
$ docker logs --since 15m --until 5m web
$ docker logs -t web
$ docker events --since 30m
$ docker events --filter type=container --filter event=die
\end{code}

Если контейнер не пишет в stdout, логи придётся искать иначе: внутри
контейнера (\cmd{docker exec}) или в привязанном каталоге на хосте.

\begin{code}[language=bash]
$ docker exec web tail -n 50 /var/log/nginx/error.log
$ sudo journalctl -u docker --since "1 hour ago"
\end{code}

\begin{notebox}
Журнал демона (\cmd{journalctl -u docker}) показывает проблемы самого
Docker: ошибки драйвера хранения, сети, запуска контейнеров. Журнал
контейнера (\cmd{docker logs})~--- проблемы приложения. Не путайте их.
\end{notebox}

\section{Ресурсы}

\begin{code}[language=bash]
$ docker stats --no-stream
$ docker stats web db
$ docker top web
$ docker system df
$ docker system df -v
\end{code}

\begin{code}[language=bash]
$ docker stats --no-stream \
    --format 'table {{.Name}}\t{{.CPUPerc}}\t{{.MemUsage}}'
NAME   CPU %   MEM USAGE / LIMIT
web    0.12%   3.4MiB / 512MiB
db     0.40%   42MiB / 1GiB
\end{code}

Если \cmd{docker stats} показывает потребление близко к лимиту,
а приложение падает с кодом 137, вероятна нехватка памяти. Проверьте
флаг \cmd{OOMKilled} и журнал ядра:

\begin{code}[language=bash]
$ docker inspect -f '{{.State.OOMKilled}}' web
true
$ sudo journalctl -k | grep -i -E 'oom|killed process' | tail
\end{code}

\section{Отладка изнутри}

\begin{code}[language=bash]
$ docker exec -it web sh
$ docker exec -u root web sh
$ docker exec -w /app web ls -la
$ docker exec web env | sort
$ docker exec web ps aux
$ docker exec web cat /etc/os-release
\end{code}

Если в образе нет оболочки и утилит (distroless, scratch), есть
несколько путей.

\subsection{Отладочный контейнер в тех же сетях и томах}

\begin{code}[language=bash]
$ docker run --rm -it \
    --network container:web \
    --pid container:web \
    nicolaka/netshoot
\end{code}

Флаг \flag{--network container:web} подключает отладочный контейнер
к сетевому пространству имён \cmd{web}, а \flag{--pid container:web}
позволяет видеть его процессы.

\subsection{docker debug}

В Docker Desktop и некоторых подписках доступна команда
\cmd{docker debug}, которая подключает к контейнеру набор отладочных
инструментов, не меняя образ.

\begin{code}[language=bash]
$ docker debug web
\end{code}

\subsection{Пространства имён хоста}

Крайняя мера~--- зайти в пространства имён контейнера с хоста.
Требует root и используется осторожно.

\begin{code}[language=bash]
$ PID=$(docker inspect -f '{{.State.Pid}}' web)
$ sudo nsenter -t "$PID" -m -u -n -i sh
\end{code}

\begin{warnbox}
Команды с \flag{--privileged}, \flag{--pid=host} и \cmd{nsenter}
дают почти полный доступ к хосту. Используйте их только на своей
машине и никогда~--- в автоматических скриптах на общих серверах.
\end{warnbox}

\section{Типичные проблемы}

\begin{center}
\small
\begin{tabularx}{\textwidth}{@{}lXX@{}}
\toprule
\textbf{Симптом} & \textbf{Причина} & \textbf{Что делать} \\
\midrule
exit 127 & команда не найдена & проверить \cmd{CMD}/\cmd{ENTRYPOINT},
\cmd{PATH}, наличие оболочки \\
exit 126 & нет прав на исполнение & \cmd{chmod +x}, проверить
архитектуру и shebang \\
exit 137 & \cmd{SIGKILL}, часто OOM & увеличить \cmd{--memory},
искать утечку, смотреть \cmd{OOMKilled} \\
exit 143 & \cmd{SIGTERM} & нормальная остановка; проверить,
кто послал сигнал \\
контейнер сразу завершается & главный процесс ушёл в фон & запускать
на переднем плане, проверить логи \\
\cmd{port is already allocated} & порт занят & \cmd{ss -ltnp},
сменить порт \\
\cmd{permission denied} на томе & несовпадение UID/GID & выровнять
владельца, \cmd{chown} в init-контейнере \\
\cmd{no space left on device} & диск заполнен & \cmd{docker system df},
prune, перенести \cmd{data-root} \\
имя не разрешается & разные сети или default bridge & общая
пользовательская сеть \\
\cmd{exec format error} & неверная архитектура & мультиарх-образ
или \flag{--platform} \\
\cmd{connection refused} & сервис слушает \cmd{127.0.0.1} & слушать
\cmd{0.0.0.0} \\
\cmd{TLS handshake timeout} & проблемы с реестром или прокси & проверить
DNS, прокси, сертификаты \\
\bottomrule
\end{tabularx}
\end{center}

\subsection{Контейнер перезапускается по кругу}

\begin{code}[language=bash]
$ docker ps -a --filter name=api
$ docker inspect -f '{{.RestartCount}}' api
$ docker logs --tail 50 api
\end{code}

Политика \cmd{always} или \cmd{unless-stopped} маскирует падение:
контейнер поднимается снова и снова. Временно отключите перезапуск
(\cmd{docker update --restart no api}), добейтесь успешного запуска
вручную и только потом возвращайте политику.

\subsection{Диск заполнен}

\begin{code}[language=bash]
$ df -h /var/lib/docker
$ sudo du -sh /var/lib/docker/* | sort -h | tail
$ docker system df -v
$ docker builder prune
$ docker image prune
$ docker container prune
\end{code}

\begin{warnbox}
\cmd{docker system prune -a --volumes} удаляет все неиспользуемые
образы, контейнеры, сети и тома. Это необратимая операция. Сначала
посмотрите \cmd{docker system df -v} и убедитесь, что среди томов
нет баз и резервных копий.
\end{warnbox}

\section{Наблюдаемость}

Для production недостаточно \cmd{docker stats}. Используйте:

\begin{itemize}
  \item \textbf{healthcheck}~--- базовый сигнал живости сервиса;
  \item \textbf{cAdvisor}~--- метрики контейнеров в формате Prometheus;
  \item \textbf{node\_exporter}~--- метрики хоста;
  \item \textbf{централизованные логи}~--- Loki, Elasticsearch,
        облачные сервисы;
  \item \textbf{трейсинг}~--- OpenTelemetry для распределённых
        запросов.
\end{itemize}

\begin{code}[language=composeyaml]
services:
  cadvisor:
    image: gcr.io/cadvisor/cadvisor:v0.49.1
    volumes:
      - /:/rootfs:ro
      - /var/run:/var/run:ro
      - /sys:/sys:ro
      - /var/lib/docker/:/var/lib/docker:ro
    ports:
      - "127.0.0.1:8080:8080"
\end{code}

\begin{notebox}
\cmd{cadvisor} требует доступа к каталогам хоста и обычно запускается
с правами, близкими к root. На shared-хостах это нужно согласовывать
с политикой безопасности.
\end{notebox}

\section{Чек-лист отладки}

\begin{enumerate}
  \item Контейнер существует и в каком состоянии?
  \item Какой код выхода и не было ли OOM?
  \item Что в логах приложения и демона?
  \item Совпадает ли команда запуска с ожидаемой?
  \item Видит ли контейнер нужные переменные, файлы и тома?
  \item В той ли сети он находится и доступен ли сосед по имени?
  \item Хватает ли ресурсов: память, CPU, диск?
  \item Не мешает ли firewall, прокси или SELinux?
\end{enumerate}

\begin{summarybox}
\begin{itemize}
  \item Начинайте с \cmd{ps}, \cmd{logs} и \cmd{inspect}.
  \item \cmd{stats}, \cmd{top} и \cmd{system df} показывают ресурсы;
        код 137 и \cmd{OOMKilled} указывают на память.
  \item Для отладки минимальных образов используйте netshoot,
        \cmd{docker debug} или \cmd{nsenter}.
  \item Политика перезапуска маскирует падения: отключайте её на время
        диагностики.
  \item В production нужны healthcheck, метрики и централизованные
        логи, а не только \cmd{docker stats}.
\end{itemize}
\end{summarybox}

\section*{Вопросы для самопроверки}
\begin{selfcheck}
  \item Какие три команды дают первичную картину состояния контейнера?
  \item Как узнать, был ли контейнер убит из-за нехватки памяти?
  \item Чем журнал демона отличается от журнала контейнера?
  \item Как отладить контейнер из distroless-образа без оболочки?
  \item Что проверять, если контейнер перезапускается по кругу?
  \item Почему \cmd{docker system prune -a --volumes} опасен?
\end{selfcheck}
